\section{Практическое применение методов верификации к L2 кэшу GPGPU}

\tab В контексте верификации нас интересуют следующие особенности кэша второго уровня: 
\begin{itemize}
    \item Наличие внутренних блоков памяти критически увеличивает количество состояний, которое может принять система
    \item Увеличенная по сравнению с кэшем первого уровня задержка предполагает увеличенное содержание последовательной логики
    \item В кэше ведется активный обмен данными, из чего следует наличие широких шин внутри системы. Это также усложняет систему.
    \item Отсутствует сложная вычислительная логика (за редким исключением)
\end{itemize}

Как уже было сказано ранее, все эти особенности не играют особой важности для методов, основанных на симуляции, поэтому в дальнейшим основной упор будет направлен на формальную верификацию, как наиболее уязвимый метод. 

Ниже будут рассматриваться отдельные блоки L2 кэша. Автор будет приводить краткое описание их функционала вместе с некоторыми требованиями из спецификации, релеватными в контексте выполнения работы. Далее будут рассматриваться разрабатываемые свойства для формальной верификации, найденные ошибки. В конце приводятся комментарии. 

\subsection{C-stage}

\subsubsection*{Краткое описание функционала}

C-stage забирает данные с D-stage, M-stage и TLB, и передает их в один из четырех выходных каналов. 

\subsubsection*{Спецификация}

\begin{itemize}
    \item входы организованы по valid-then-ready handshake. 
    \item недопустим одновременный прием с более чем одного входа 
    \item Выбор между D и M интерфейсами происходит по Round Robin принципу
    \item TLB более приоритетный, чем D и M 
    \item Входные пакеты распределяются одному из четырех интерфейсов в соответствии с номером, содержащимся в пакете
    \item Сигнал o\_pkt\_cnt отвечает за число пакетов, содержащихся в очереди соответствующего интерфейса
\end{itemize}

\subsubsection*{Верификация}

Рукопожатие верифицируем по аналогии. 

Для проверки арбитрации введем секвенции, соответствующие проведенной транзакции на каждом интерфейсе. Тогда, можем записать группу свойств:

Особую сложность для формальной верификации представляют счетчики. В нашем случае он небольшой, и можно было бы его верифицировать с помощью другого, референсного счетчика, но гораздо экономнее будет проверить условия инкремента, декремента и отсутствия изменений.

Для проверки корректной передачи данных используем свойства следующего типа:

\subsubsection*{Найденные ошибки и их анализ}

Неправильная арбитрация: при определенных условиях мог быть выбран клиент с более низкой приоритетностью  % d10039b5

Найдена ошибка в bypass buffer: сигнал, что буффер полон, не сохранялся когда данные не отправлялись в или из буффера % f43b2951

\subsection{D-stage}

\subsubsection*{Кратное описание функционала}

D-stage забирает инструкции с данными с S-stage, при необходимости записывает, читает и/или выполняет ATOMIC инструкцию, а после отправляет выходной пакет на C, M, Q-stage.

\subsubsection*{Спецификация}

\begin{itemize}
    \item valid-then-ready handshake
    \item Операция MEMBAR просто проходит дальше на C интерфейс
    \item Операция STORE пишет данные в память, на выход ничего не подает. 
    \item Операция LOAD читает данные из памяти и пишет их в C интерфейс.
    \item Операция ATOM читает данные из памяти, отправляет запрос во Atomic Unit, отправляет считанные данные в C интерфейс, после чего записывает в память результат выполненной операции
\end{itemize}

\subsubsection*{Верификация}

Основаную сложность в верификации этого модуля представляет его объем и вычислительная нагрузка, связанная с ATOMIC операциями. Без каких либо упрощений, даже простая компиляция вызыывает у инструмента трудности, поэтому для начала по очередии разберемся с этими прблемами. 

Во-первых, попробуем решить проблему наличия памяти. Блок SRAM 512*2048 бит находится за гранью любого совресенного инструмента формальной верификации. Сам по себе он создает 2\^1067008 состояний, не включая сигналов доступа к нему. Можно было бы попробовать условно разделить блок на две части: до и после памяти. В первой части мы будем проверять, что в память записываются правильные величины, подключаясь непосредственно к управляющим сигналам; во второй части будет проверять правильность выходных данных, поключаясь к сигналам из памяти. 

Второй подход заключается в замещении реально памяти ее моделью. JasperGold предоставляет для некоторых часто используемых, но сложных для формальной верификации, уствройств специальные блоки, являющиеся упрощенными моделями данных устройств. Такие блоки будет называть акселераторами, согласно документации. Сам акселератор работает по следующему принципу: внутри себя он может хранить до 32 величин, соответствеющих паре адрес+данные. Т.к. в процессе верификации блоков с памятью внутри одного свойства редко требуется больше одного-двух адресов, то такая упрощенная модель функционально ничем не отличается от реальной памяти, при этом сильно ускоряя процесс доказательства.

Самый главный минус акселератора --- его крайняя степень негибкости. Сигналы, контролирующие запись и чтение, могут работать отлично от реализации в проверяемой модели, что требует дополниьельных модификаций. В данному случае модификации оказались минимальными, тем не менее, подключение акселератора к модели заняло суещственную часть времени разработки свойств, что заметно отложило получение первых результатов. 

Второй важной проблемой является начилие внутри логически простого блока устроства для проведения атомик операция, в т.ч. над числами с плавающей запятой. Известно, что формальная верификация не подходит дял верификации таких блоков, поэтому, ввиду отсутисвия других решений, было принято решение подключчаться только к входным и выходным сигналам модуля, не првоеряя его внутренню логику. Таким образов, мы исключали ее из конуса влияния, чем еще сильнее упрощаем верицикацию.

Для valid-then-ready рукопожатия напишем свойтсва следующего вида:

\begin{lstlisting}
property p_rdy;
    !i_vld |-> o_rdy
    endproperty
\end{lstlisting}

\subsubsection*{Найденные ошибки и их анализ}

Нерпавильная посылка данных в ATOM unit: Операция ATOM пишет свой результат в неправильный адрес, если новая инструкция пришла (т.е. поднялся ipktvld) раньше, чем закончил выполняться ATOM. Например, ATOM сменяется LOAD, который читает из адреса отличного от того, куда пишет ATOM. Сигнал opktrdy равен 0, т.е LOAD еще не начал обрабатываться, однако адрес, куда записывается результат ATOM, меняется. % 2d9d56a4

Если пакет, обращающийся к L2 Atomic, несколько тактов не сменяется другим пакетом, то D-stage непрерывно обращается к L2 Atomic. % ae34985d

% При определенных условиях атом операции фризились % 8713a216

Проверка на ipktvld=1 при STORE/UPDATE отсутствует и может произойти запись данных из невалидного пакета % 662в3388

\subsection{DD-stage}

\subsubsection*{Кратное описание функционала}

DD-stage 

\subsubsection*{Спецификация}

Направляет TLBREQ в TLB
Направляет STORE/ATOM/RED/MEMBAR/FLUSH/INVAL в T
Делит LOAD на 1-4 LOADа и направляет в T
Делает PRIC приходящего адреса, если это не FLUSH, INVAL, TLBREQ, после сериализации LOAD
Команды FLUSH/INVAL приводят к отправке 2048 пакетов FLUSH/EVICT в T. Новые пакеты не принимаются, пока не пришло 2048 (строго) ответов от T/S/Q.

\begin{itemize}
    \item 
\end{itemize}

\subsubsection*{Верификация}

Верификация счетчика выполняется аналогично предыдущему блоку. Единственное отличие --- теперь счетчик может одновременно инкрементироваться на 1, 2 или 3. Дл упрощения, разделим на 3 свойства, каждое из которых будет проверять инкремент на соответствующее число: 

Особую сложность представляет верификация деления операции LOAD. Количество операций зависит от флага внутри пришедшего LOAD. Для поднятия сигнала ordy после отправки всех LOAD напишем следующее свойство:

Проверка сериализации данных выполняется несколько сложнее. В исходном пакете присутствует 4 поля для максимально возможного количества выходных пакетов, в каждом из полей записаны данные, соответствующие конкретному выходному пакету. Необоходимо проверить запись данных, учитывая то, что пакетов до рассматриваемого может быть любое количество. Для этого запишем следующее свойство:

PRIC адресы выполняется с помощью функции

Остальное тривиально

\subsubsection*{Найденные ошибки и их анализ}

Найдено не было

\subsection{T-stage}

\subsubsection*{Кратное описание функционала}

 Управляет наличием cache-линии в кэше. И запиранием CacheLine для атомиков. PLRU

\subsubsection*{Спецификация}

\begin{itemize}
    \item 
\end{itemize}

\subsubsection*{Верификация}

T-stage можно условно разделить на PLRU, T0, T1 и память. Верифицировать блок целиком нецелесообразно --- в нем происходит множество логических преобразований, так что нетривиальное свойство, использующее только входные и выходных сигналы, будет истекать по таймауту компиляции. Другим неочевидным преимущество такого разделения блока на подблоки является локализация ошибки. 

Особое внимание уделим алгоритму PLRU.

Для вериифкации компаратора входного тэга с тэгами из памяти: 

Для раздвоения операций:



\subsubsection*{Найденные ошибки и их анализ}

Последняя транзакция (evict), должна была передать на выход oevicttag, который был у предыдущего store(way == 0), но передает oevicttag с way == 8 % 

\subsection{Q-stage}

\subsubsection*{Кратное описание функционала}

Буферизирует Pending-запросы в Memory X-bar в reqq, отправляет запросы в память и кладёт их в очередь ожидания ответа: misslookupq, flushlookupq, тем самым предотвращает дэдлок L2Partition. Также буферизирует прочитанные данные в updateq очереди. Может отправить флаг о переполнении в T-stage.

\subsubsection*{Спецификация}

\begin{itemize}
    \item 
\end{itemize}

\subsubsection*{Верификация}



\subsubsection*{Найденные ошибки и их анализ}

Не обнаружено 









